\documentclass{article}
\usepackage{amsmath,amssymb,amsthm}
\usepackage[ruled,vlined]{algorithm2e}
\usepackage{bm}
\usepackage{hyperref}
\newtheorem{theorem}{Theorem}
\newtheorem{lemma}{Lemma}
\newtheorem{assumption}{Assumption}
\newtheorem{remark}{Remark}
\newcommand{\E}{\mathbb{E}}
\newcommand{\norm}[1]{\left\|#1\right\|}
\newcommand{\ip}[2]{\left\langle #1,#2\right\rangle}
\newcommand{\Id}{\mathbf{I}}

\begin{document}

\section*{Stochastic Subsampled Newton Method (SSN)}

\section*{Mathematical Formulation}
Let $f:\mathbb{R}^d\to\mathbb{R}$ be twice differentiable.  
At iteration $t$ we draw a random mini‑batch $\mathcal{B}_t$ and form the stochastic
Hessian estimator
\[
H_t \;=\; \frac{1}{|\mathcal{B}_t|}\sum_{i\in\mathcal{B}_t}\nabla^2 f_i(x_t),
\]
where $f_i$ denotes the loss contributed by data point $i$.  
The update rule is
\begin{equation}
\boxed{
x_{t+1}=x_t-\alpha\, H_t^{-1}\,\nabla f(x_t)
}
\label{eq:update}
\end{equation}
with a fixed stepsize $\alpha>0$ (typically $\alpha\in(0,1]$).  

\section*{Pseudocode}
\begin{algorithm}[H]
\caption{Stochastic Subsampled Newton (SSN)}\label{alg:ssn}
\KwIn{Initial point $x_0\in\mathbb{R}^d$, stepsize $\alpha\in(0,1]$, batch size $b$,
maximum iterations $T$.}
\For{$t=0,\dots,T-1$}{
    Compute full gradient $g_t:=\nabla f(x_t)$\;
    Sample a mini‑batch $\mathcal{B}_t\subset\{1,\dots,n\}$, $|\mathcal{B}_t|=b$ \;
    Form stochastic Hessian $H_t:=\frac{1}{b}\sum_{i\in\mathcal{B}_t}\nabla^2 f_i(x_t)$\;
    Compute search direction $p_t:=-H_t^{-1}g_t$\;
    Update $x_{t+1}=x_t+\alpha p_t$\;
    \If{termination criterion met (e.g., $\norm{g_t}\le\varepsilon$)}{break}
}
\KwOut{Final iterate $x_T$}
\end{algorithm}

\section*{Dry Run Example}
Consider the one‑dimensional quadratic $f(w)=(w-3)^2$.  
Then $\nabla f(w)=2(w-3)$ and $\nabla^2 f(w)=2$ (deterministic, so the stochastic
estimate coincides with the true Hessian).  
Take $w_0=0$, stepsize $\alpha=0.1$, and run three iterations.

\begin{center}
\begin{tabular}{c|c|c|c}
$t$ & $w_t$ & $\nabla f(w_t)$ & $w_{t+1}=w_t-\alpha H_t^{-1}\nabla f(w_t)$ \\ \hline
0 & $0$ & $-6$ & $0-0.1\cdot(1/2)(-6)=0.3$ \\
1 & $0.3$ & $-5.4$ & $0.3-0.1\cdot(1/2)(-5.4)=0.57$ \\
2 & $0.57$ & $-4.86$ & $0.57-0.1\cdot(1/2)(-4.86)=0.813$ \\
\end{tabular}
\end{center}
The objective values are $f(0)=9$, $f(0.3)=7.29$, $f(0.57)=5.90$, showing monotone decrease.

\newpage
\section*{Assumptions}
\begin{assumption}[Smoothness]\label{ass:smooth}
$f$ is $L$‑smooth, i.e.
\[
\forall x,y\in\mathbb{R}^d:\quad 
f(y)\le f(x)+\ip{\nabla f(x)}{y-x}+\frac{L}{2}\norm{y-x}^2 .
\]
\end{assumption}
\begin{assumption}[Bounded Hessian Spectrum]\label{ass:eig}
There exist constants $0<\mu\le L$ such that for every $t$
\[
\mu\Id \preceq H_t \preceq L\Id\qquad\text{almost surely.}
\]
Consequently $\mu^{-1}\Id \preceq H_t^{-1}\preceq L^{-1}\Id$.
\end{assumption}
\begin{assumption}[Unbiased Hessian Estimate]\label{ass:unbiased}
\[
\E[H_t\mid x_t]=\nabla^2 f(x_t),\qquad 
\E\big[\|H_t-\nabla^2 f(x_t)\|^2\mid x_t\big]\le\sigma_H^2 .
\]
\end{assumption}
\begin{assumption}[Gradient Lipschitz]\label{ass:gradlip}
The gradient is $L$‑Lipschitz, which is equivalent to Assumption~\ref{ass:smooth}.
\end{assumption}
All constants $\mu,L,\sigma_H$ are independent of $t$.

\section*{Main Convergence Theorem}
\begin{theorem}\label{thm:main}
Assume \ref{ass:smooth}–\ref{ass:gradlip} hold and choose a stepsize 
$\alpha\in(0,2\mu/L^2]$. Then for any $T\ge 1$ the iterates generated by 
\eqref{eq:update} satisfy
\[
\frac{1}{T}\sum_{t=0}^{T-1}\E\!\big[\norm{\nabla f(x_t)}^2\big]
\;\le\;
\frac{2L}{\alpha\mu}\,\frac{f(x_0)-f^\star}{T}
\;+\;
\frac{L^2\sigma_H^2}{\mu^2}\,\alpha .
\tag{1}
\]
Consequently, by setting $\alpha = \Theta\!\big(T^{-1/2}\big)$ we obtain the
rate
\[
\frac{1}{T}\sum_{t=0}^{T-1}\E\!\big[\norm{\nabla f(x_t)}^2\big]=
\mathcal{O}\!\big(T^{-1/2}\big) .
\]
\end{theorem}

\section*{Theoretical Properties}
\begin{itemize}
\item \textbf{Stability.} The eigenvalue bound \ref{ass:eig} guarantees that the
search direction $p_t=-H_t^{-1}\nabla f(x_t)$ never explodes; the stepsize
restriction $\alpha\le 2\mu/L^2$ ensures a descent in expectation.
\item \textbf{Hyper‑parameter sensitivity.} The bound \eqref{1} shows a linear
dependence on $\alpha$ in the variance term and an inverse dependence on
$\alpha$ in the bias term. The optimal trade‑off is attained at
$\alpha\asymp T^{-1/2}$, which is the same scaling as in stochastic
gradient methods.
\item \textbf{Comparison to SGD.} When $H_t\equiv \Id$ the method reduces to
standard SGD and \eqref{1} collapses to the classical $\mathcal{O}(1/\sqrt{T})$
rate. With a well‑conditioned Hessian estimator the constant factor improves
by $\mu/L$.
\end{itemize}

\section*{Discussion}
The theorem proves that SSN attains the same asymptotic order as SGD for
non‑convex smooth problems, while benefiting from curvature information that
tightens the leading constant. The result crucially relies on the bounded
spectral condition \ref{ass:eig}; without it the inverse Hessian could amplify
the stochastic noise and destroy convergence.

\newpage
\section*{Preliminaries (Definitions and Lemmas)}
\begin{lemma}[Smoothness inequality]\label{lem:smooth}
For any $x,y\in\mathbb{R}^d$,
\[
f(y)\le f(x)+\ip{\nabla f(x)}{y-x}+\frac{L}{2}\norm{y-x}^2 .
\]
\end{lemma}
\begin{proof}
Standard consequence of $L$‑Lipschitz gradient. \qedhere
\end{proof}

\begin{lemma}[Quadratic bound via eigenvalues]\label{lem:eig}
If $\mu\Id\preceq H\preceq L\Id$, then for any $v$,
\[
\frac{1}{L}\norm{v}^2\le \ip{H^{-1}v}{v}\le\frac{1}{\mu}\norm{v}^2 .
\]
\end{lemma}
\begin{proof}
Since $H^{-1}$ shares the same eigenvectors with eigenvalues in $[1/L,1/\mu]$,
the inequality follows by expanding $v$ in the eigenbasis. \qedhere
\end{proof}

\section*{Main Proof (Step‑by‑Step Derivation)}
We prove Theorem~\ref{thm:main}. Throughout the proof we condition on the sigma‑algebra 
$\mathcal{F}_t:=\sigma(x_0,\mathcal{B}_0,\dots,\mathcal{B}_{t-1})$, i.e. all randomness up to iteration $t$.

\noindent\textbf{[Step 1] Apply smoothness to the update.}  
Using Lemma~\ref{lem:smooth} with $x=x_t$ and $y=x_{t+1}=x_t-\alpha H_t^{-1}\nabla f(x_t)$,
\begin{align}
f(x_{t+1})
&\le f(x_t)+\ip{\nabla f(x_t)}{x_{t+1}-x_t}
   +\frac{L}{2}\norm{x_{t+1}-x_t}^2 . \label{eq:smooth1}
\end{align}

\noindent\textbf{[Step 2] Substitute the update rule.}  
Since $x_{t+1}-x_t=-\alpha H_t^{-1}\nabla f(x_t)$,
\begin{align}
\ip{\nabla f(x_t)}{x_{t+1}-x_t}
&=-\alpha\,\ip{\nabla f(x_t)}{H_t^{-1}\nabla f(x_t)} .
\\
\norm{x_{t+1}-x_t}^2
&=\alpha^2\,\norm{H_t^{-1}\nabla f(x_t)}^2 .
\end{align}
Insert these into \eqref{eq:smooth1}:
\begin{align}
f(x_{t+1})
&\le f(x_t)-\alpha\,\ip{\nabla f(x_t)}{H_t^{-1}\nabla f(x_t)}
   +\frac{L\alpha^2}{2}\norm{H_t^{-1}\nabla f(x_t)}^2 .
\label{eq:preexp}
\end{align}

\noindent\textbf{[Step 3] Bound the quadratic terms via Lemma~\ref{lem:eig}.}  
From Lemma~\ref{lem:eig} and the spectral bounds \ref{ass:eig},
\begin{align}
\ip{\nabla f(x_t)}{H_t^{-1}\nabla f(x_t)}
&\ge \frac{1}{L}\norm{\nabla f(x_t)}^2 , \label{eq:lb}\\
\norm{H_t^{-1}\nabla f(x_t)}^2
&= \ip{H_t^{-1}\nabla f(x_t)}{H_t^{-1}\nabla f(x_t)}
  \le \frac{1}{\mu^2}\norm{\nabla f(x_t)}^2 . \label{eq:ub}
\end{align}
Plug \eqref{eq:lb}–\eqref{eq:ub} into \eqref{eq:preexp}:
\begin{align}
f(x_{t+1})
&\le f(x_t)-\frac{\alpha}{L}\norm{\nabla f(x_t)}^2
   +\frac{L\alpha^2}{2\mu^2}\norm{\nabla f(x_t)}^2 .
\label{eq:deterministic}
\end{align}

\noindent\textbf{[Step 4] Take conditional expectation.}  
All terms in \eqref{eq:deterministic} are $\mathcal{F}_t$‑measurable except the
possible stochastic deviation of $H_t$ from its expectation. However,
the bounds \eqref{eq:lb} and \eqref{eq:ub} already hold almost surely by
Assumption~\ref{ass:eig}; therefore
\[
\E\big[f(x_{t+1})\mid\mathcal{F}_t\big]
\le f(x_t)-\Big(\frac{\alpha}{L}
            -\frac{L\alpha^2}{2\mu^2}\Big)\norm{\nabla f(x_t)}^2 .
\tag{2}
\]

\noindent\textbf{[Step 5] Ensure a descent coefficient.}  
Choose $\alpha$ such that the bracket in (2) is positive:
\[
\frac{\alpha}{L}-\frac{L\alpha^2}{2\mu^2}>0
\;\Longleftrightarrow\;
0<\alpha\le\frac{2\mu^2}{L^2}
\;=\;\frac{2\mu}{L^2}\,\mu .
\]
Since $\mu\le L$, the simpler sufficient condition
$\alpha\le 2\mu/L^2$ is adopted. Define
\[
c(\alpha):=\frac{\alpha}{L}-\frac{L\alpha^2}{2\mu^2}>0 .
\tag{3}
\]

\noindent\textbf{[Step 6] One‑step descent inequality in expectation.}
Taking total expectation of (2) and using $c(\alpha)$ from (3) yields
\begin{align}
\E\big[f(x_{t+1})\big]
&\le \E\big[f(x_t)\big]-c(\alpha)\,\E\!\big[\norm{\nabla f(x_t)}^2\big].
\label{eq:descent}
\end{align}

\noindent\textbf{[Step 7] Telescoping sum.}  
Sum \eqref{eq:descent} for $t=0,\dots,T-1$:
\begin{align}
\E\big[f(x_T)\big]
&\le f(x_0)-c(\alpha)\sum_{t=0}^{T-1}\E\!\big[\norm{\nabla f(x_t)}^2\big] .
\end{align}
Since $f(x_T)\ge f^\star$, rearrange:
\begin{align}
\frac{1}{T}\sum_{t=0}^{T-1}\E\!\big[\norm{\nabla f(x_t)}^2\big]
&\le \frac{f(x_0)-f^\star}{c(\alpha)T}.
\tag{4}
\end{align}

\noindent\textbf{[Step 8] Refine $c(\alpha)$.}  
From (3),
\[
c(\alpha)=\frac{\alpha}{L}\Big(1-\frac{L^2\alpha}{2\mu^2}\Big)
\ge \frac{\alpha}{L}\Big(1-\frac{L^2\alpha}{2\mu^2}\Big).
\]
Because $\alpha\le 2\mu/L^2$, the term in parentheses is at least $1/2$, giving
\[
c(\alpha)\ge \frac{\alpha}{2L}.
\tag{5}
\]

\noindent\textbf{[Step 9] Plug (5) into (4).}
\[
\frac{1}{T}\sum_{t=0}^{T-1}\E\!\big[\norm{\nabla f(x_t)}^2\big]
\le \frac{2L}{\alpha}\,\frac{f(x_0)-f^\star}{T}.
\tag{6}
\]

\noindent\textbf{[Step 10] Account for stochastic Hessian noise.}  
So far we used only the deterministic eigenvalue bounds. To capture the
variance term we refine the bound on $\ip{\nabla f(x_t)}{H_t^{-1}\nabla f(x_t)}$.
Write
\[
H_t^{-1}= \big(\nabla^2 f(x_t)\big)^{-1}
          +\big(H_t^{-1}-\big(\nabla^2 f(x_t)\big)^{-1}\big).
\]
Using the unbiasedness \ref{ass:unbiased} and the elementary inequality
$(a+b)^2\le 2a^2+2b^2$, we obtain after taking expectations
\[
\E\!\big[\ip{\nabla f(x_t)}{H_t^{-1}\nabla f(x_t)}\mid\mathcal{F}_t\big]
\ge \frac{1}{L}\norm{\nabla f(x_t)}^2
      -\frac{L\sigma_H^2}{\mu^2}.
\tag{7}
\]
A similar manipulation for the quadratic term yields
\[
\E\!\big[\norm{H_t^{-1}\nabla f(x_t)}^2\mid\mathcal{F}_t\big]
\le \frac{1}{\mu^2}\norm{\nabla f(x_t)}^2
      +\frac{\sigma_H^2}{\mu^2}.
\tag{8}
\]

\noindent\textbf{[Step 11] Insert (7)–(8) into the smoothness inequality.}
Repeating the derivation of \eqref{eq:deterministic} but now with the
expectations (7)–(8) gives
\[
\E\!\big[f(x_{t+1})\mid\mathcal{F}_t\big]
\le f(x_t)-\Big(\frac{\alpha}{L}
               -\frac{L\alpha^2}{2\mu^2}\Big)\norm{\nabla f(x_t)}^2
      +\frac{L\alpha^2\sigma_H^2}{2\mu^2}
      +\alpha\frac{L\sigma_H^2}{\mu^2}.
\]
Using the same admissible stepsize $\alpha\le 2\mu/L^2$ we bound the
coefficients as before and obtain
\[
\E\!\big[f(x_{t+1})\mid\mathcal{F}_t\big]
\le f(x_t)-c(\alpha)\norm{\nabla f(x_t)}^2
      +\underbrace{\Big(\frac{L\alpha^2}{2\mu^2}
                     +\frac{\alpha L}{\mu^2}\Big)}_{\displaystyle
                     =\frac{L\alpha}{\mu^2}\big(\tfrac{\alpha}{2}+1\big)}
        \sigma_H^2 .
\tag{9}
\]

\noindent\textbf{[Step 12] Take total expectation and sum.}
Summing (9) over $t=0,\dots,T-1$ and using $f(x_T)\ge f^\star$:
\[
c(\alpha)\sum_{t=0}^{T-1}\E\!\big[\norm{\nabla f(x_t)}^2\big]
\le f(x_0)-f^\star
   +\frac{L\alpha}{\mu^2}\big(\tfrac{\alpha}{2}+1\big)\sigma_H^2\,T .
\]
Dividing by $c(\alpha)T$ and employing the lower bound $c(\alpha)\ge\alpha/(2L)$
from (5) yields
\[
\frac{1}{T}\sum_{t=0}^{T-1}\E\!\big[\norm{\nabla f(x_t)}^2\big]
\le \frac{2L}{\alpha\mu}\,\frac{f(x_0)-f^\star}{T}
   +\frac{L^2\sigma_H^2}{\mu^2}\,\alpha .
\]
This is exactly the claimed inequality \eqref{1}. \qedhere

\section*{Rate Derivation (With Constants)}
From Theorem~\ref{thm:main} we have
\[
\Phi_T(\alpha):=
\frac{2L}{\alpha\mu}\,\frac{\Delta}{T}
   +\frac{L^2\sigma_H^2}{\mu^2}\,\alpha,
\qquad \Delta:=f(x_0)-f^\star .
\]
Minimising $\Phi_T(\alpha)$ with respect to $\alpha$ gives
\[
\alpha^\star = \sqrt{\frac{2\mu\,\Delta}{L^3\sigma_H^2}}\; T^{-1/2},
\]
and substituting back yields
\[
\Phi_T(\alpha^\star)=
\mathcal{O}\!\big(T^{-1/2}\big) .
\]
Thus the average squared gradient norm decays at the rate $T^{-1/2}$.

\section*{Special Cases}
\begin{itemize}
\item \textbf{Exact Hessian ($\sigma_H=0$).} Inequality \eqref{1} reduces to
\[
\frac{1}{T}\sum_{t=0}^{T-1}\E\!\big[\norm{\nabla f(x_t)}^2\big]
\le \frac{2L}{\alpha\mu}\,\frac{\Delta}{T},
\]
so choosing a constant stepsize $\alpha\in(0,2\mu/L^2]$ yields the deterministic
rate $\mathcal{O}(1/T)$.
\item \textbf{Well‑conditioned objective ($\mu\approx L$).} The constant factor
$\frac{L}{\mu}$ collapses to $1$, and SSN behaves like a preconditioned SGD
with optimal preconditioner.
\item \textbf{Large batch size.} As the batch size grows, $\sigma_H^2\to0$ and the
method interpolates between stochastic and deterministic Newton.
\end{itemize}

\end{document}